\section*{Вопрос 1}

Строгого определения ИИ нет. В учебнике Павлова перечислено больше десятка вариантов, и каждый цепляет свою сторону вопроса. Сначала разведём два термина. Искусственный интеллект --- это область исследований и свойство поведения машины, а система искусственного интеллекта --- конкретная программа. Мы строим как раз программу, поэтому важнее второе.

Из определений я считаю рабочими три.

Первое --- у И. Рич и Дж. Алена. ИИ это область, где создают компьютеры, выполняющие то, что пока лучше получается у человека. Определение честное. Задача раньше требовала человека и его опыт, теперь её решает программа. Минус в том, что про устройство системы оно ничего не говорит.

Второе --- определение СИИ у Павлова. Это компьютерная креативная система со сложной структурой, которая накапливает и корректирует знания, ставит и достигает цель и при этом адаптируется к изменениям среды и собственного состояния. Здесь уже названы проверяемые свойства. Знания, цель, адаптация. Там же есть краткий вариант. СИИ это аппаратно-программный комплекс, способный накапливать и корректировать знания и действовать целенаправленно на основе логического вывода.

Третье --- формула из лекции Хохловой. Обычная программа это алгоритм плюс данные. СИИ это знания плюс стратегия обработки знаний. Коротко и по сути. Главное отличие в том, что система работает со знаниями, которые лежат отдельно от кода и могут пополняться. Там же названы четыре признака знаний. Внутренняя интерпретируемость, структурированность, связанность и активность.

Своё определение я бы собрал из этих трёх так. Система искусственного интеллекта --- это программа, которая хранит знания о предметной области отдельно от механизма их обработки, умеет их пополнять и на их основе вырабатывать решение для цели, которую раньше достигал человек. Такое определение легко проверить на любой программе. Есть ли знания, отделены ли они от кода, пополняются ли по ходу работы, есть ли цель. Если да, систему можно считать интеллектуальной по делу, а не по названию.
